\section{Additional Background}
\label{appendix:sec:additional_background}

%We provide more details on background here, and we note that they are introduced in a way that facilitates discussing the proposed method and that other works~\cite{tfn, 3dsteerable, kondor2018clebsch, cormorant, se3_transformer, segnn} also provide relevant mathematical background, which helps the presentation of this work.   

%%%In this section, we provide additional mathematical background on group equivariance helpful for the discussion of the proposed method.
In this section, we provide additional mathematical background helpful for the discussion of the proposed method.
Other works~\citep{tfn, 3dsteerable, kondor2018clebsch, cormorant, se3_transformer, segnn} also provide similar background.
We encourage interested readers to see these works~\citep{zee, Dresselhaus2007} for more in-depth and pedagogical presentations.

\subsection{Group Theory}
\label{appendix:subsec:group_theory}

\paragraph{Definition of Groups.}
A group is an algebraic structure that consists of a set $G$ and a binary operator $\circ: G \times G \rightarrow G$ and is typically denoted as $G$.
Groups satisfy the following four axioms:
\begin{enumerate}
    \item Closure: $g \circ h \in G$ for all $g, h \in G$.
    \item Identity: There exists an identity element $e \in G$ such that $g \circ e = e \circ g = g$ for all $g \in G$.
    \item Inverse: For each $g \in G$, there exists an inverse element $g^{-1} \in G$ such that $g \circ g^{-1} = g^{-1} \circ g = e$.
    \item Associativity: $f \circ g \circ h = (f \circ g) \circ h = f \circ (g \circ h )$ for all $f, g, h \in G$.
\end{enumerate}

In this work, we focus on 3D rotation, translation and inversion.
Relevant groups include:
\begin{enumerate}
    \item The Euclidean group in three dimensions $E(3)$: 3D rotation, translation and inversion.
    \item The special Euclidean group in three dimensions $SE(3)$: 3D rotation and translation.
    \item The orthogonal group in three dimensions  $O(3)$: 3D rotation and inversion.
    \item The special orthogonal group in three dimensions  $SO(3)$: 3D rotation.
    %\item The translational group $T(3)$: 3D translation.
\end{enumerate}


\paragraph{Group Representations.}
%Groups can be used to define transformations.
The actions of groups define transformations.
Formally, a transformation acting on vector space $X$ parametrized by group element $g \in G$ is an injective function $T_{g}: X \rightarrow X$.
A powerful result of group representation theory is that these transformations can be expressed as matrices which act on vector spaces via matrix multiplication. 
These matrices are called the group representations.
Formally, a group representation $D: G \rightarrow GL(N)$ is a mapping between a group $G$ and a set of $N \times N$ invertible matrices.
The group representation $D(g): X \rightarrow X$ maps an $N$-dimensional vector space $X$ onto itself and satisfies  $D(g)D(h) = D(g \circ h)$ for all $g, h \in G$.

%There are many ways to represent a group. 
%If there exists a change of basis $P$ in the form of an $N \times N$ matrix for a given representation $D(g)$ such that $P^{-1} D(g) P = D'(g)$
How a group is represented depends on the vector space it acts on.
If there exists a change of basis $P$ in the form of an $N \times N$ matrix such that $P^{-1} D(g) P = D'(g)$ for all $g \in G$, then we say the two group representations are equivalent.
If $D'(g)$ is block diagonal, which means that $g$ acts on independent subspaces of the vector space, the representation $D(g)$ is reducible.
%that makes $D'(g)$ block diagonal for all $g \in G$, meaning $g$ acts on independent subspaces of the vector space, the representation $D(g)$ is deemed reducible. 
%A particular class of representations that are convenient for generalizable and composable functions are irreducible representations or ``irreps'', which cannot be further reduced. 
A particular class of representations that are convenient for composable functions are irreducible representations or ``irreps'', which cannot be further reduced. 
We can express any group representation of $SO(3)$ as a direct sum (concatentation) of irreps \citep{zee,Dresselhaus2007,e3nn}:
\begin{equation}
D(g) = P ^{-1} \left ( \bigoplus_{i} D_{l_i}(g) \right)  P = 
P^{-1}
\begin{pmatrix}
D_{l_0}(g) & & \\
& D_{l_1}(g) & & \\
& & ...... \\
\end{pmatrix}
P
\end{equation}
where $D_{l_i}(g)$ are Wigner-D matrices with degree $l_i$ as metnioned in Sec.~\ref{subsec:irreducible_representations}.

% already mention in background
%\paragraph{Irreducible Representations of $SO(3)$.}



\subsection{Equivariance}
\label{appendix:subsec:equivariance}


\paragraph{Definition of Equivariance and Invariance.}
Equivariance is a property of a function $f: X \rightarrow Y$ mapping between vector spaces $X$ and $Y$.
Given a group $G$ and group representations $D_{X}(g)$ and $D_{Y}(g)$ in input and output spaces $X$ and $Y$, $f$ is equivariant to G if $D_Y(g) f(x) = f(D_X(g)x)$ for all $x \in X$ and $g \in G$.
%Invariance corresponds to the case that $f(x) = D_Y(g) f(x) = f(D_X(g) x)$.
Invariance corresponds to the case where $D_Y(g)$ is the identity $I$ for all $g \in G$.


\paragraph{Equivariance in Neural Networks.}
%Including equivariance in neural networks can enable generalizaing to unseen data in a predictable manner.
%For instance, an equivariant neural network $f: X \rightarrow Y$ predicts output $y$ given input $x$.
%When the input becomes $D_X(g) x$, which might not exist in training sets, $f$ can generalize to the input and predict $D_Y(g) f(x)$.
Group equivariant neural networks are guaranteed to to make equivariant predictions on data transformed by a group.
Additionally, they are found to be data-efficient and generalize better than non-symmetry-aware and invariant methods~\citep{nequip, quantum_scaling, neural_scale_of_chemical_models}.
For 3D atomistic graphs, we consider equivariance to the Euclidean group $E(3)$, which consists of 3D rotation, translation and inversion.
For translation, we operate on relative positions and therefore our networks are invariant to 3D translation.
%As for rotation and inversion, we incorporate vector spaces of $O(3)$ irreps into networks so that intermediate features and outputs are equivariant to 3D rotation and inversion.
We achieve equivariance to rotation and inversion by representing our input data, intermediate features and outputs in vector spaces of $O(3)$ irreps and acting on them with only equivariant operations. 


\subsection{Equivariant Features Based on Vector Spaces of Irreducible Representations}
\label{appendix:subsec:equivariant_irreps_feature}


\paragraph{Irreducible Representations of Inversion.}
The group of inversion $\mathbb{Z}_2$ only has two elements, identity and inversion, and two irreps, even $e$ and odd $o$.
Vectors transformed by irrep $e$ do not change sign under inversion while those by irrep $o$ do.
We create irreps of $O(3)$ by simply multiplying those of $SO(3)$ and $\mathbb{Z}_2$ and introduce parity $p$ to type-$L$ vectors to denote how they transform under inversion.
%Therefore for $O(3)$, type-$L$ vectors in $SO(3)$ are extended to type-$(L, p)$ vectors, where $p$ is $e$ or $o$.
Thus, type-$L$ vectors in $SO(3)$ become type-$(L, p)$ vectors in $O(3)$, where $p$ is $e$ or $o$.
%In the following, we use type-$L$ vectors for the ease of discussion, but we can generalize to type-$(L, p)$ vectors, unless otherwise stated.}

\paragraph{Irreps Features.}
As discussed in Sec.~\ref{subsec:irreducible_representations} in the main text, we use type-$L$ vectors for $SE(3)$-equivariant irreps features\footnote{In SEGNN~\citep{segnn}, they are also referred to as steerable features. We use the term ``irreps features'' to remain consistent with \texttt{e3nn}~\citep{e3nn} library.} and type-$(L, p)$ vectors for $E(3)$-equivariant irreps features.
Parity $p$ denotes whether vectors change sign under inversion and can be either $e$ (even) or $o$ (odd).
Vectors with $p = o$ change sign under inversion while those with $p = e$ do not.
Scalar features correspond to type-$0$ vectors in the case of $SE(3)$-equivariance and correspond to type-$(0, e)$ in the case of $E(3)$-equivariance whereas type-$(0, o)$ vectors correspond to pseudo-scalars.
Euclidean vectors in $\mathbb{R}^3$ correspond to type-$1$ vectors and type-$(1, o)$ vectors whereas type-$(1, e)$ vectors correspond to pseudo-vectors.
Note that type-$(L, e)$ vectors and type-$(L, o)$ vectors are considered vectors of different types in equivariant linear layers and layer normalizations.

\paragraph{Spherical Harmonics.}
Euclidean vectors $\vec{r}$ in $\mathbb{R}^3$ can be projected into type-$L$ vectors $f^{(L)}$ by using spherical harmonics $Y^{(L)}$: $f^{(L)} = Y^{(L)}(\frac{\vec{r}}{|| \vec{r} ||})$~\citep{finding_symmetry_breaking_order}.
This is equivalent to the Fourier transform of the angular degree of freedom $\frac{\vec{r}}{|| \vec{r}||}$, which can be optionally weighted by $|| \vec{r} ||$. 
In the case of $SE(3)$-equivariance, $f^{(L)}$ transforms in the same manner as type-$L$ vectors.
For $E(3)$-equivariance, $f^{(L)}$ behaves as type-$(L, p)$ vectors, where $p = e$ if $L$ is even and $p = o$ if $L$ is odd.
\revision{
Visualization of spherical harmonics can be found in this \href{https://docs.e3nn.org/en/latest/api/o3/o3_sh.html}{website}.}

%\paragraph{\todo{Spherical Harmonics are Wigner-D matrices.}}

\paragraph{Vectors of Higher $L$ and Other Parities.}
%So far we have restricted concrete examples of $O(3)$ irreps to commonly encountered scalars ($L=0$, $p=1$) and vectors ($L=1$, $p=-1$); irreps of higher $L$'s and other parities are equally physical. 
Although previously we have restricted concrete examples of vector spaces of $O(3)$ irreps to commonly encountered scalars (type-$(0, e)$ vectors) and Euclidean vectors (type-$(1, o)$ vectors), vector of higher $L$ and other parities are equally physical. 
%Many physical properties are described using Cartesian tensors
For example, the moment of inertia (how an object rotates under torque) transforms as a $3\times3$ symmetric matrix, which has symmetric-traceless components behaving as type-$(2, e)$ vectors.
Elasticity (how an object deforms under loading) transforms as a rank-4 or $3\times3\times3\times3$ symmetric tensor, which includes components acting as type-$(4, e)$ vectors.
%As an example, many physical properties are described using Cartesian tensors, e.g. the moment of inertia (how an object rotates under torque) transforms as a $3\times3$ symmetric matrix and elasticity (how an object deforms under loading) transforms as a rank-4 or $3\times3\times3\times3$ symmetric tensor, which have components that transform as ($L=2$, $p=1$) and ($L=4$, $p=1$), respectively. 
%\todo{More generally speaking, if one considers polynomials of distinguishable 3D vectors $\vec{x}^{(i)}$, specific linear combinations of monomials will transform as specific irreps of $O(3)$. In fact, the spherical harmonics are a specific subset of these polynomials where the inputs are identical $\vec{x}^(i) \rightarrow \vec{x}$ and the degree of the polynomial matches the $L$ of the irrep (e.g. $xy$ is of degree 2 and transforms as a component of ($L=2$, $p=1$)) \cite{e3nn}.}


\subsection{Tensor Product}
\label{appendix:subsec:tensor_product}

\paragraph{Tensor Product for $O(3)$.}
We use tensor products to interact different type-$(L, p)$ vectors.
We extend our discussion in Sec.~\ref{subsec:tensor_product} in the main text to include inversion and type-$(L, p)$ vectors.
The tensor product denoted as $\otimes$ uses Clebsch-Gordan coefficients to combine type-$(L_1, p_1)$ vector $f^{(L_1, p_1)}$ and type-$(L_2, p_2)$ vector $g^{(L_2, p_2)}$ and produces type-$(L_3, p_3)$ vector $h^{(L_3, p_3)}$ as follows:
\begin{equation}
%\begin{aligned}
h^{(L_3, p_3)}_{m_3} = (f^{(L_1, p_1)} \otimes g^{(L_2, p_2)})_{m_3} = \sum_{m_1 = -L_1}^{L_1} \sum_{m_2 = -L_2}^{L_2} C^{(L_3, m_3)}_{(L_1, m_1)(L_2, m_2)} f^{(L_1, p_1)}_{m_1} g^{(L_2, p_2)}_{m_2} 
\label{appendix:eq:tensor_product}
\end{equation}

\begin{equation}
p_3 = p_1 \times p_2
\label{appendix:eq:parity}
\end{equation}

The only difference of tensor products for $O(3)$ as described in Eq.~\ref{appendix:eq:tensor_product} from those for $SO(3)$ described in Eq.~\ref{eq:tensor_product} is that we additionally keep track of the output parity $p_3$ as in Eq.~\ref{appendix:eq:parity} and use the following multiplication rules: $e \times e = e$, $o \times o = e$, and $e \times o = o \times e = o$.
For example, the tensor product of a type-$(1, o)$ vector and a type-$(1, e)$ vector can result in one type-$(0, o)$ vector, one type-$(1, o)$ vector, and one type-$(2, o)$ vector.


\paragraph{Clebsch-Gordan Coefficients.}
The Clebsch-Gordan coefficients for $SO(3)$ are computed from integrals over the basis functions of a given irreducible representation, e.g., the real spherical harmonics, as shown below and are tabulated to avoid unnecessary computation.
\begin{equation}
     C_{(L_1, m_1) (L_2, m_2)}^{(L_3, m_3)} = \left|L_1 m_1; L_2 m_2 \right> \left<L_3 m_3 \right| = \int d{\Omega} Y^{(L_1)*}_{m_1}(\Omega) Y^{(L_2)*}_{m_2}(\Omega) Y^{(L_3)}_{m_3}(\Omega)
\end{equation}
For many combinations of $L_1$, $L_2$, and $L_3$, the Clebsch-Gordan coefficients are zero. The gives rise to the following selection rule for non-trivial coefficients: $-|L_1 + L_2| \leq L_3 \leq |L_1 + L_2|$.


\paragraph{Examples of Tensor Products.}
Tensor products generally define the interaction between different type-$(L, p)$ vectors in a symmetry-preserving manner and consist of common operations as follows:
\begin{enumerate}
    \item Scalar-scalar multiplication: scalar $(L=0, p = e)$ $\otimes$ scalar $(L=0, p=e)$ $\rightarrow$ scalar $(L=0, p=e)$.
    \item Scalar-vector multiplication: scalar $(L=0, p = e)$ $\otimes$ vector $(L=1, p=o)$ $\rightarrow$ vector $(L=1, p=o)$.
    \item Vector dot product: vector $(L=1, p = o)$ $\otimes$ vector $(L=1, p=o)$ $\rightarrow$ scalar $(L=0, p=e)$.
    \item Vector cross product: vector $(L=1, p = o)$ $\otimes$ vector $(L=1, p=o)$ $\rightarrow$ pseudo-vector $(L=1, p=e)$.
\end{enumerate}

%\subsection{$E(3)$-Equivariant Neural Networks}
%Euclidean neural networks are built such that ever operation is equivariant to Euclidean symmetry, satisfying the condition for $g \in E(3)$, $f:X \times W \rightarrow Y$, inputs $x$ and scalar weights $w$, $f(D_X(g)x, w) = D_Y(g)f(x, w)$ \cite{tfn,kondor2018clebsch,3dsteerable}. While these networks have primarily been used within the framework of 3D convolutional neural networks \cite{3dsteerable} and message-passing graph neural networks \cite{tfn,kondor2018clebsch,nequip} the mathematical framework is more general and can be used for other modalities \cite{allergo,se3_transformer}. We use the $E(3)$ equivariant, parameterizable components in Ref. \cite{e3nn} for building our networks.